\begin{proposition}[The normalization and modification algorithm]\hypertarget{tag-02ZN}{}\label{push:src:A.2.5}\label{push:A25}
\StatementRecord{push:src:A.2.5}{proved}{02ZN}{\SeriesRef{push:src:A.2.4}, \SeriesRef{push:src:A.4.6}, \SeriesRef{push:src:A.4.8}, \SeriesRef{push:src:A.2.2}}

For every finite fan datum $\Sigma$ of Theorem~\ref{push:src:A.2.4}, the following algorithm describes the normalized components and their comparison with a reference model.

The star charts and integer transition matrices of \ref{push:src:A.2.4}
give every normalized component, its gluing maps, and
all multiple loci directly. There is also a terminating
finite-data algorithm that expresses its comparison
with a reference model as an explicit sequence of
blow-ups and blow-downs along specified smooth strata.
At each blow-up the algorithm gives the normalized
strict transforms, the exceptional projective bundle,
its multiplicity, and every new intersection.
The algorithm uses both blow-ups and blow-downs.

For a component of multiplicity $\ell$, the local
nonreduced central equation is the monomial in \eqref{pub:I:eq:A-24d}.
More generally, in a unimodular chart with central
equation $F=\prod_i z_i^{\ell_i}$, the cotangent-complex
first-order deformation sheaf is
\begin{equation}
 \mathcal T^1_D
  =\mathcal O_D\Big/
    \left(\ell_i z_i^{\ell_i-1}
          \prod_{j\ne i}z_j^{\ell_j}:\ell_i>0\right).
\label{pub:I:eq:A-25a}
\end{equation}
Here $\mathcal T^1_D$ means
$\mathcal Ext^1(\mathbb L_D,\mathcal O_D)$; in the
reduced case it also equals
$\mathcal Ext^1(\Omega_D,\mathcal O_D)$.
The reduced case is the trivial line on the singular
scheme, as in \ref{push:src:A.2.4}. For a nonreduced centre the displayed
Jacobian quotient is the appropriate deformation sheaf;
the stated d-semistability conclusion applies to the reduced case.
\end{proposition}


\begin{proof}
\begin{enumerate}[label=\arabic*.,leftmargin=*]
\item Proposition~\ref{push:src:A.2.3} proves that the point-star component $\mathbb P^3$ has no proper bimeromorphic map to $R$. In that example the full normalization is $\operatorname{Bl}_pR^+\sqcup\mathbb P^3$: the first component is the strict transform, and the two meet along the exceptional plane. These are the ordinary point-blow-up charts of the smooth total fourfold.

\item Retain all marking
data from \ref{push:src:A.2.4}: the lattice, the two period matrices,
the original line-bundle units, and the common
unsubdivided fan. Strictify the finite quotient cone
complex by the decreasing-dimensional barycentric
stars used in \ref{push:src:A.4.8}. Record these operations and do
the same for the reference model. Form their common
intersection complex. A regular common subdivision
can be found either by toric desingularization or by
enumerating rational subdivisions with increasing
primitive-ray bounds and testing their determinants
and common faces. A regular subdivision exists, so
the latter enumeration terminates.

\item On a finite strict regular complex there are finitely
many possible ordinary star blow-ups: choose a cone
of dimension at least two and insert the sum of its
primitive basis rays. There are also finitely many
possible inverse moves: choose a ray and test whether
its star is exactly the star created by such a
subdivision of a regular cone. Both tests are exact
integer calculations. Require that every intermediate
complex remain a subdivision of the marked original
complex. Record face identifications; equality of two
marked complexes is a finite rational comparison,
using the period orbit representatives in \eqref{pub:I:eq:A-24c} and
their pure top and pure bottom analogues.

\item Enumerate this locally finite graph of moves in
increasing path length, discarding repeated marked
complexes. Applied separately to the two proper
subdivision morphisms from the common regular
complex, the finite-complex weak-factorization theorem
used in \ref{push:src:A.4.8} guarantees a path, all of whose
intermediates are over the corresponding target.
Thus the search terminates. Concatenate the two paths
and the recorded strictification paths. This gives a
fully specified finite sequence once the input fan is
given. Temporary pure bottom subdivisions may occur;
they are allowed in a presentation of the birational
comparison, and are removed before the prescribed
endpoint is reached.

\item  If a cone has primitive basis vectors
$v_1,\ldots,v_c$, its centre is the closed stratum
\[
 C=\{z_1=\cdots=z_c=0\}
\]
in the marked toroidal chart. Strictness makes its
global closure smooth; the period translates defining
one centre are disjoint. Blow-up chart $i$ has
coordinates
\[
 z_i=t_i,\qquad z_j=t_i w_j\quad
 (1\leq j\leq c,\ j\ne i),
 \qquad z_j\quad(j>c).
\]
The changes between charts are these integral
monomials, with the original units raised to the
same powers. The exceptional divisor is
$\mathbb P_C(N_{C/Y})$, where $Y$ is the smooth
intermediate total space. If $C$ is central, its
central multiplicity is
\[
 \ell_E=\sum_{j=1}^c\ell_j.
\]
If all $\ell_j$ vanish, it is a horizontal exceptional
divisor and is not a new central component.

\item For an old normalized component, take the inverse
image of this centre on its marked branch and the
ordinary strict transform in these charts. It is
the blow-up along that smooth inverse-image centre;
when that centre is Cartier the blow-up is the
identity. Its intersections with the new exceptional
component are the appropriate coordinate-hyperplane
projective subbundles. Intersections of several old
branches are given by imposing the corresponding
coordinate equations simultaneously. Thus the same
chart prescription records double, triple and
quadruple loci, their normal lines, and their 
gluing maps. An inverse star is the inverse
contraction of this recorded configuration. This
specifies the centres and their order.

\item The algorithm starts with the explicit reference
normalizations of \ref{push:src:A.2.2} and the $n$-cycle reference
in \ref{push:src:A.4.8}. It is therefore a uniform geometric
presentation of every output by specified blow-ups,
blow-downs and projective bundles. Independently,
the direct star charts avoid running this search
when only a normalization chart or an intersection
is required. The mixed-ray components are complete
smooth toric threefolds; a pure top component
$(0,b,0,\ell)$ is bimeromorphic to
$\mathbb P_S(\mathcal O\oplus L^\ell)$.

\item Finally $D$ is globally the principal Cartier
divisor of the parameter $r$ in a smooth total
space. Its two-term cotangent complex is
$[\mathcal O_D\xrightarrow{dF}\Omega_Y|_D]$,
so its first-order deformation module is the
cokernel of the gradient. Equivalently, a lifted
hypersurface equation $F+\epsilon G$ is changed by
an infinitesimal coordinate change precisely by a
Jacobian term. This argument applies to nonreduced
monomials as well; an ordinary conormal sequence
need not be left-exact there. Absorbing the local unit, or reducing its
derivative modulo $F$, gives \eqref{pub:I:eq:A-25a}. A global
parameter makes the normal line trivial, so these
Jacobian quotients glue by the canonical
hypersurface presentation. In the reduced case
the Jacobian ideal defines the singular scheme,
giving the asserted trivial deformation line there.

\end{enumerate}
\end{proof}
